% Este documento foi adaptado por Marcos Laureano a partir do original disponibilizado
% em https://www.overleaf.com/latex/templates/modelo-canonico-do-ifpr-com-logo/mqtrmrvzyfsb feito por Rafael Silva Santos

\documentclass[12pt, openright, oneside, a4paper, english, brazilian]{abntex2}


% ---
% Pacotes
% ---
\usepackage{lmodern}			% Usa a fonte Latin Modern			
\usepackage[T1]{fontenc}		% Selecao de codigos de fonte.
\usepackage[utf8]{inputenc}		% Codificacao do documento (conversão automática dos acentos)
\usepackage{indentfirst}		% Indenta o primeiro parágrafo de cada seção.
\usepackage{color}				% Controle das cores
\usepackage{graphicx}			% Inclusão de gráficos
\usepackage{microtype} 			% para melhorias de justificação
\usepackage[brazilian,hyperpageref]{backref}	 % Paginas com as citações na bibl
\usepackage[alf]{abntex2cite}	% Citações padrão ABNT
\usepackage{lipsum}     %Gerador de texto aleatorio em Latim

\usepackage[normalem]{ulem} %sublinhar linhas e permitir quebrar linhas
\usepackage{amssymb} %necessário para $\square$


% --- 
% CONFIGURAÇÕES DE PACOTES
% --- 

% ---
% Configurações do pacote backref
% Usado sem a opção hyperpageref de backref
\renewcommand{\backrefpagesname}{Citado na(s) página(s):~}
% Texto padrão antes do número das páginas
\renewcommand{\backref}{}
% Define os textos da citação
\renewcommand*{\backrefalt}[4]{
	\ifcase #1 %
		Nenhuma citação no texto.%
	\or
		Citado na página #2.%
	\else
		Citado #1 vezes nas páginas #2.%
	\fi}%
% ---



% ----------------------------------
% --- Configurações da capa do IFPR
% ----------------------------------
\renewcommand{\imprimircapa}{%
    \begin{capa}%
        \center
        \begin{figure*}
            \centering
            \includegraphics[height=2.5cm]{figuras/logo-ifpr.png}
        \end{figure*}
        \vspace*{\fill}
        {\ABNTEXchapterfont\large\imprimirautor}\\
        \vspace*{\fill}
        \vspace{2.5cm}{\ABNTEXchapterfont\bfseries\LARGE\imprimirtitulo}\\
        \vspace*{\fill}
        \vspace*{\fill}
        \vspace{2.5cm}
        {\large\imprimirlocal}\\
        \par
        {\large\imprimirdata}\\
        \vspace*{1cm}
    \end{capa}
}


% ---
% Informações de dados para CAPA e FOLHA DE ROSTO
% ---
\titulo{GUIA PRÁTICO DE INFRAESTRUTURA EM NUVEM: Passo a passo para hospedar um servidor de \textit{Minecraft} de alta performance gratuito}    %Título

\autor{Gabriel de Paula Gaspar Pinto}  %Autores
\local{Pinhais}
\data{2026}
% \orientador{Prof. Fulano de Tal}
% \coorientador{Prof. Deltrano de Tal}
\instituicao{%
  Instituto Federal do Paraná -- IFPR
  \par
  Campus Pinhais
  \par
  Bacharelado em Ciência da Computação}
\tipotrabalho{Relatório de Atividades de Extensão}       % Tipo de trabalho Monografia de Conclusão de Curso, Tese de Doutorado, Dissertação de Mestrado etc.


% O preambulo deve conter o tipo do trabalho, o objetivo, o nome da instituição e a área de concentração. Preambulo aparece na folha de Rosto
\preambulo{Relatório de atividades extensionistas, apresentado à disciplina de Práticas de Extensão III, do curso de Bacharelado em Ciência da Computação do Instituto Federal do Paraná, Campus Pinhais, sob a orientação do professor Marcos Aurelio Pchek Laureano}
% ---


% ---
% Configurações de aparência do PDF final

% alterando o aspecto da cor azul
\definecolor{black}{RGB}{0,0,0}




% informações do PDF
\makeatletter
\hypersetup{
     	%pagebackref=true,
		pdftitle={\@title}, 
		pdfauthor={\@author},
    	pdfsubject={\imprimirpreambulo},
	    pdfcreator={LaTeX with abnTeX2},
		pdfkeywords={abnt}{latex}{abntex}{abntex2}{trabalho acadêmico}, 
		colorlinks=true,       		% false: boxed links; true: colored links
    	linkcolor=black,          	% color of internal links
    	citecolor=black,        		% color of links to bibliography
    	filecolor=magenta,      		% color of file links
		urlcolor=black,
		bookmarksdepth=4
}
\makeatother
% --- 

% ---
% Posiciona figuras e tabelas no topo da página quando adicionadas sozinhas
% em um página em branco. Ver https://github.com/abntex/abntex2/issues/170
\makeatletter
\setlength{\@fptop}{5pt} % Set distance from top of page to first float
\makeatother
% ---

% ---
% Possibilita criação de Quadros e Lista de quadros.
% Ver https://github.com/abntex/abntex2/issues/176
%
\newcommand{\quadroname}{Quadro}
\newcommand{\listofquadrosname}{Lista de quadros}

\newfloat[chapter]{quadro}{loq}{\quadroname}
\newlistof{listofquadros}{loq}{\listofquadrosname}
\newlistentry{quadro}{loq}{0}

% configurações para atender às regras da ABNT
\setfloatadjustment{quadro}{\centering}
\counterwithout{quadro}{chapter}
\renewcommand{\cftquadroname}{\quadroname\space} 
\renewcommand*{\cftquadroaftersnum}{\hfill--\hfill}

\setfloatlocations{quadro}{hbtp} % Ver https://github.com/abntex/abntex2/issues/176
% ---

% --- 
% Espaçamentos entre linhas e parágrafos 
% --- 

% O tamanho do parágrafo é dado por:
\setlength{\parindent}{1.3cm}

% Controle do espaçamento entre um parágrafo e outro:
\setlength{\parskip}{0.2cm}  % tente também \onelineskip

% ---
% compila o indice
% ---
\makeindex
% ---

% ----
% Início do documento
% ----
 %% aqui você deve alterar dados como nomes dos autores e orientadores

\begin{document}


% Seleciona o idioma do documento (conforme pacotes do babel)
%\selectlanguage{english}
\selectlanguage{brazilian}

% Retira espaço extra obsoleto entre as frases.
\frenchspacing 

% ----------------------------------------------------------
% ELEMENTOS PRÉ-TEXTUAIS
% ----------------------------------------------------------
% \pretextual

% ---
% Capa OBRIGATÓRIO
% ---
\imprimircapa
% ---

% ---
% Folha de rosto OBRIGATÓRIO
% (o * indica que haverá a ficha bibliográfica)
% ---
\imprimirfolhaderosto*
% ---

% ---
% Inserir a ficha bibliografica OPCIONAL
% ---

% Isto é um exemplo de Ficha Catalográfica, ou ``Dados internacionais de
% catalogação-na-publicação''. Você pode utilizar este modelo como referência. 
% Porém, provavelmente a biblioteca da sua universidade lhe fornecerá um PDF
% com a ficha catalográfica definitiva após a defesa do trabalho. Quando estiver
% com o documento, salve-o como PDF no diretório do seu projeto e substitua todo
% o conteúdo de implementação deste arquivo pelo comando abaixo:


\begin{fichacatalografica}
	\sffamily
	\vspace*{\fill}					% Posição vertical
	\begin{center}					% Minipage Centralizado
	\fbox{\begin{minipage}[c][8cm]{13.5cm}		% Largura
	\small
	\imprimirautor
	%Sobrenome, Nome do autor
	
	\hspace{0.5cm} \imprimirtitulo  / \imprimirautor. --
	\imprimirlocal, \imprimirdata-
	
	\hspace{0.5cm} \thelastpage p. : il. (algumas color.) ; 30 cm.\\
	
	\hspace{0.5cm} \imprimirorientadorRotulo~\imprimirorientador\\
	
	\hspace{0.5cm}
	\parbox[t]{\textwidth}{\imprimirtipotrabalho~--~\imprimirinstituicao,
	\imprimirdata.}\\
	
	\hspace{0.5cm}
		1. Palavra-chave1.
		2. Palavra-chave2.
		2. Palavra-chave3.
		I. Orientador.
		II. Instituto Federal do Paraná - Campus Pinhais.
		III. Título 			
	\end{minipage}}
	\end{center}
\end{fichacatalografica}
% ---

% ---
% Folha de aprovação. OBRIGATÓRIO em trabalhos de conclusão de curso. OPCIONAL nos demais trabalhos acadêmicos.
% ---

% Isto é um exemplo de Folha de aprovação, elemento obrigatório da NBR
% 14724/2011 (seção 4.2.1.3). Você pode utilizar este modelo até a aprovação
% do trabalho. Após isso, substitua todo o conteúdo deste arquivo por uma
% imagem da página assinada pela banca com o comando abaixo:
%
% \begin{folhadeaprovacao}
% \includepdf{folhadeaprovacao_final.pdf}
% \end{folhadeaprovacao}
%
\begin{folhadeaprovacao}

  \begin{center}
    {\ABNTEXchapterfont\large\imprimirautor}

    \vspace*{\fill}\vspace*{\fill}
    \begin{center}
      \ABNTEXchapterfont\bfseries\Large\imprimirtitulo
    \end{center}
    \vspace*{\fill}
    
    \hspace{.45\textwidth}
    \begin{minipage}{.5\textwidth}
        \imprimirpreambulo
    \end{minipage}%
    \vspace*{\fill}
   \end{center}
        
   Trabalho aprovado. \imprimirlocal, dd de Mês de yyyy:

   \assinatura{\textbf{\imprimirorientador} \\ Orientador} 
   \assinatura{\textbf{Professor} \\ Convidado 1}
   \assinatura{\textbf{Professor} \\ Convidado 2}
   %\assinatura{\textbf{Professor} \\ Convidado 3}
   %\assinatura{\textbf{Professor} \\ Convidado 4}
      
   \begin{center}
    \vspace*{0.5cm}
    {\large\imprimirlocal}
    \par
    {\large\imprimirdata}
    \vspace*{1cm}
  \end{center}
  
\end{folhadeaprovacao}
% ---
 %%aqui também tem coisas que precisam ser alteradas
% % ---
% % Dedicatória OPCIONAL
% % ---
% % \begin{dedicatoria}
% %    \vspace*{\fill}
% %    \centering
% %    \noindent
% %    \textit{Elemento Opcional. Este trabalho é dedicado às crianças adultas que,\\
% %    quando pequenas, sonharam em se tornar cientistas.} \vspace*{\fill}
% % \end{dedicatoria}
% % ---

% % ---
% % Agradecimentos OPCIONAL
% % ---
% \begin{agradecimentos}
% Os agradecimentos principais são direcionados à Gerald Weber, Miguel Frasson,
% Leslie H. Watter, Bruno Parente Lima, Flávio de Vasconcellos Corrêa, Otavio Real
% Salvador, Renato Machnievscz\footnote{Os nomes dos integrantes do primeiro
% projeto abn\TeX\ foram extraídos de
% \url{http://codigolivre.org.br/projects/abntex/}} e todos aqueles que
% contribuíram para que a produção de trabalhos acadêmicos conforme
% as normas ABNT com \LaTeX\ fosse possível.

% Agradecimentos especiais são direcionados ao Centro de Pesquisa em Arquitetura
% da Informação\footnote{\url{http://www.cpai.unb.br/}} da Universidade de
% Brasília (CPAI), ao grupo de usuários
% \emph{latex-br}\footnote{\url{http://groups.google.com/group/latex-br}} e aos
% novos voluntários do grupo
% \emph{\abnTeX}\footnote{\url{http://groups.google.com/group/abntex2} e
% \url{http://www.abntex.net.br/}}~que contribuíram e que ainda
% contribuirão para a evolução do \abnTeX.

% \end{agradecimentos}
% % ---

% % ---
% % Epígrafe OPCIONAL
% % ---
% \begin{epigrafe}
%     \vspace*{\fill}
% 	\begin{flushright}
% 		\textit{``Não vos amoldeis às estruturas deste mundo, \\
% 		mas transformai-vos pela renovação da mente, \\
% 		a fim de distinguir qual é a vontade de Deus: \\
% 		o que é bom, o que Lhe é agradável, o que é perfeito.\\
% 		(Bíblia Sagrada, Romanos 12, 2)}
% 	\end{flushright}
% \end{epigrafe}
% % ---


% ---
% RESUMOS
% ---

% resumo em português OBRIGATÓRIO
\setlength{\absparsep}{18pt} % ajusta o espaçamento dos parágrafos do resumo
\begin{resumo}
 Descreva aqui o resumo do seu trabalho.
 
 \textbf{Palavras-chave}: Computação em Nuvem. Administração de Sistemas. Servidores Virtuais Privados (VPS). Jogos Digitais. Guia Prático.  
\end{resumo}

% resumo em inglês OBRIGATÓRIO
\begin{resumo}[Abstract]
 \begin{otherlanguage*}{english}
   O mesmo, só que em inglês.

   \vspace{\onelineskip}
 
   \noindent 
   \textbf{Keywords}: Cloud Computing. System Administration. Virtual Private Servers (VPS). Video Games. Practical Guide.
 \end{otherlanguage*}
\end{resumo}


% ---
% inserir lista de ilustrações OPCIONAL
% ---
\pdfbookmark[0]{\listfigurename}{lof}
\listoffigures*
\cleardoublepage
% ---

% ---
% inserir lista de quadros OPCIONAL
% ---
\pdfbookmark[0]{\listofquadrosname}{loq}
\listofquadros*
\cleardoublepage
% ---

% ---
% inserir lista de tabelas OPCIONAL
% ---
\pdfbookmark[0]{\listtablename}{lot}
\listoftables*
\cleardoublepage
% ---

% ---
% inserir lista de abreviaturas e siglas OPCIONAL
% ---
\begin{siglas}
  \item[ABNT] Associação Brasileira de Normas Técnicas
  \item[abnTeX] ABsurdas Normas para TeX
\end{siglas}
% ---

% ---
% inserir lista de símbolos OPCIONAL
% ---
\begin{simbolos}
  \item[$ \Gamma $] Letra grega Gama
  \item[$ \Lambda $] Lambda
  \item[$ \zeta $] Letra grega minúscula zeta
  \item[$ \in $] Pertence
\end{simbolos}
% ---




% ---
% inserir o sumario OBRIGATÓRIO
% ---
\pdfbookmark[0]{\contentsname}{toc}
\tableofcontents*
\cleardoublepage

% ---
% ----------------------------------------------------------
% ELEMENTOS TEXTUAIS
% ----------------------------------------------------------
\textual
%%%%%%%%%%%%%%%%%%%%%%%%%%%%%%%%%%%%%%%%%%%%%%%%%%%%%%%%%%%%%%%
%% TEXTO AQUI
%%%%%%%%%%%%%%%%%%%%%%%%%%%%%%%%%%%%%%%%%%%%%%%%%%%%%%%%%%%%%%%


\chapter{INTRODUÇÃO}
\label{capitulo-introducao}
    A convergência entre práticas recreativas no ambiente digital e a aplicação de conceitos fundamentais de Ciência da Computação constitui uma maneira de extensão para difusão científica e tecnológica. Dentre as plataformas digitais interativas, o jogo eletrônico \textit{Minecraft} se destaca não apenas por sua amplitude global, mas também por sua arquitetura cliente-servidor descentralizada, que atua como porta de entrada para que usuários desenvolvam interesses em redes, administração de servidores e protocolos de comunicação.

    No contexto acadêmico do curso de Bacharelado em Ciência da Computação do Instituto Federal do Paraná (IFPR) - Câmpus Pinhais, para a disciplina de Práticas de Extensão III, esta ação foi estruturada para atender a uma demanda de jogadores, entusiastas de tecnologia e comunidades locais, a incapacidade técnica e financeira de manter servidores persistentes, de alta performance e seguros, sem incorrer em custos recorrentes, incompatíveis com a realidade da sociedade brasileira.

    A proposta será concretizada na elaboração, gravação e disponibilização de uma série de videoaulas práticas em formato de playlist pública no YouTube, cujo objetivo central é ensinar o provisionamento e a manutenção autônoma de um ambiente virtual de computação em nuvem. A metodologia adotada descomplica o uso de uma infraestrutura corporativa ao demonstrar como utilizar a camada de gratuidade permanente (\textit{Always Free}) da Oracle Cloud Infrastructure (OCI). O escopo de instruções abrange desde a criação da instância virtual até práticas recomendadas de segurança de borda, incluindo o roteamento de portas com \textit{iptables} e listas de controle de acesso (\textit{Security Lists}), gerenciamento de zonas DNS e túneis criptografados com a Cloudflare, filtragem geográfica de pacotes e a publicação de servidores web com Nginx.

    Dessa forma, o projeto articula as diretrizes da extensão universitária ao transferir competências práticas e acadêmicas de engenharia de software e infraestrutura de redes diretamente para a sociedade civil, promovendo a capacitação autônoma de administradores e transformando o lazer em experiência educacional aplicada.


\chapter{DELIMITAÇÃO DO PROBLEMA}
\label{capitulo-delimitacao}

	Neste momento, procure destacar o problema de seu trabalho. Aborde de maneira objetiva a pergunta-problema que instiga a sua investigação.

	Introduza questões reflexivas, tais como:
	\begin{itemize}
		\item Que questões-chave precisam ser respondidas?
		\item Quais as implicações práticas da ação de extensão? 
		\item Quais foram as principais lacunas nesta proposta de extensão?
		\item Quais foram os principais desafios nesta proposta de trabalho?
	\end{itemize}


\chapter{JUSTIFICATIVA}
\label{capitulo-justificativa}
A realização desta atividade de extensão justifica-se pela oportunidade de aproximar conceitos de computação de ponta da realidade cotidiana da comunidade. Frequentemente, a infraestrutura em nuvem corporativa e a administração de sistemas baseados em Linux são vistas como temas complexos e de difícil absorção para pessoas fora do ambiente profissional de TI. Utilizar o ecossistema do \textit{Minecraft} como estudo de caso prático funciona como um estímulo de engajamento, permitindo que conceitos mais complexos de redes de computadores, roteamento e segurança sejam apreendidos a partir de um projeto simples e de utilidade imediata.

Sob a perspectiva econômica e técnica, o projeto valida uma alternativa de custo zero, cuja capacidade supera as soluções comerciais populares de entrada. Enquanto opções pagas no mercado brasileiro exigem mensalidades recorrentes por instâncias restritas a pouca memória e processadores limitados, a camada \textit{Always Free} da Oracle Cloud oferece, na arquitetura Ampere ARM, até 2 OCPUs e 12 GB de RAM. Com essas configurações disponibilizadas e a vinculação de um endereço IPv4 público, com latência baixa, devido à proximidade do local onde o servidor está localizado (São Paulo/SP), permitindo uma autonomia operacional integral.

Além da economia financeira direta para os usuários atendidos, a relevância da proposta também está na abrangência dos tópicos técnicos ensinados. O conhecimento adquirido não se restringe à operação do jogo. O participante é capacitado para administrar um nó Linux completo, habilitando-o a instalar aplicações corporativas ou de uso pessoal, como centrais de mídia (\textit{Plex ou Jellyfin}), servidores de correio eletrônico, servidores web de documentação ou até mesmo túneis de rede privada.

Se a ação não fosse concretizada, os usuários permaneceriam sujeitos a custos contínuos desnecessários ou restritos a ambientes de baixa confiabilidade, comprovando a visão de que a computação em nuvem é inacessível para fins não corporativos.


\chapter{OBJETIVO}
\label{capitulo-objetivo}
	Apresente, aqui, os seus objetivos com o desenvolvimento da ação de extensão, distinguindo o objetivo geral dos objetivos específicos. Descreva o objetivo da atividade prática realizada, de forma clara e sucinta. 
	Busque iniciar por um verbo \textit{dicendi}, por exemplo:
	\begin{itemize}
		\item Estudar a aplicação de…
		\item Descrever… 
		\item Implementar… 
		\item Identificar… 
		\item Implantar… 
		\item Levantar/ gerar dados… 
		\item Apresentar/ elaborar… 
		\item Promover…
		\item Fomentar…
	\end{itemize}

\chapter{METODOLOGIA}
\label{capitulo-metodologia}
    A metodologia deste trabalho estrutura-se a partir de duas frentes que se complementam, sendo elas o desenvolvimento e a consolidação de uma infraestrutura em nuvem funcional e o processo pedagógico de transposição didática desse conhecimento em um produto educacional em formato de videoaulas.

	\section{Materiais}
    Para a concepção do ambiente computacional e a produção do material didático, foram empregados os recursos de infraestrutura, softwares de rede e utilitários descritos a seguir:

\subsection{Recursos de Nuvem e Hardware}
\begin{itemize}
    \item \textbf{Oracle Cloud Infrastructure (OCI) - Camada \textit{Always Free}:} Instância de computação baseada em arquitetura ARM (\textit{Ampere}, dimensionada com 2 OCPUs, 12 GB de memória RAM e volume de armazenamento em bloco (\textit{Boot Volume}) de 200GB, associada a um endereço IPv4 público atribuído dinamicamente que permanece fixo enquanto a instância estiver ativa.
    \item \textbf{Local de Trabalho:} Computador pessoal utilizado para acesso administrativo remoto via terminal, geração e envio de chaves criptográficas SSH.
\end{itemize}

\subsection{Sistemas Operacionais, Servidores e Plataformas de Rede}
\begin{itemize}
    \item \textbf{Sistema Operacional Hospedeiro:} Distribuição Linux (\textit{Ubuntu Server}) executada na instância da OCI.
    \item \textbf{Núcleo de Aplicação de Jogos:} \textit{PaperMC}, configurado em conjunto com a \textit{Java Virtual Machine} (JVM) parametrizada pelas \textit{Aikar's flags} para otimização da coleta de lixo (\textit{Garbage Collection}).
    \item \textbf{Módulos de Compatibilidade e Interface:} \textit{GeyserMC} e \textit{Floodgate} para viabilizar o acesso simultâneo entre clientes das edições Java e Bedrock; \textit{BlueMap} para a renderização de mapas tridimensionais do mundo gerado; e \textit{MiniMOTD} para padronização estética da lista de servidores.
    \item \textbf{Servidor Web Reverso:} \textit{Nginx}, empregado na recepção e encaminhamento do tráfego HTTP/HTTPS direcionado às interfaces web do servidor.
    \item \textbf{Serviços de DNS e Conectividade de Borda:} Provedor \textit{DuckDNS} para resolução dinâmica de nomes na etapa inicial; plataforma \textit{Registro.br} para aquisição e delegação de autoridade do domínio de topo (\textit{Top-Level Domain} - TLD); e \textit{Cloudflare} para o gerenciamento autoritativo de zonas DNS, mitigação de ataques.
    \item \textbf{Túneis Criptografados e Segurança de Borda:} Daemon \textit{cloudflared} para criação do \textit{Cloudflare Tunnel} dedicado ao tráfego do protocolo SSH, listas de controle de rede (\textit{Security Lists} na OCI) e \textit{iptables} para filtragem local de pacotes.
\end{itemize}


	\section{Métodos}
	\label{capitulo-metodologia-metodos}
    O método adotado classifica-se como uma pesquisa aplicada e de caráter descritivo, fundamentada na documentação empírica de um caso de uso real de engenharia de infraestrutura. O fluxo metodológico foi dividido em cinco etapas sequenciais:

\subsection{Etapa 1: Provisionamento e Parametrização da Camada de Rede Inicial}
A primeira fase consistiu na criação da conta na OCI e no provisionamento de uma máquina virtual com processador ARM Ampere, configurada com 2 OCPUs, 12 GB de RAM e sistema operacional Linux. Com a máquina alocada e o endereço IPv4 público associado, estabeleceram-se as primeiras regras de comunicação:
\begin{enumerate}
    \item Liberação seletiva nas \textit{Security Lists} da Virtual Cloud Network (VCN) da Oracle para permitir tráfego de entrada (\textit{ingress}) nas portas essenciais de administração e operação;
    \item Adequação e persistência de regras de filtragem em nível de núcleo com o \textit{iptables}, tratando o redirecionamento de portas e evitando bloqueios involuntários originados das diretivas padrão da imagem do sistema operacional.
\end{enumerate}

\subsection{Etapa 2: Resolução de Nomes Dinâmica e Implantação da Aplicação}
Com a infraestrutura de rede acessível, procedeu-se à instalação dos componentes de software:
\begin{enumerate}
    \item Configuração de um serviço de DNS Dinâmico (DDNS) através do \textit{DuckDNS}, viabilizando um endereço nominal gratuito para testes preliminares de conexão sem dependência direta do endereço IP numérico;
    \item Instalação do ambiente de execução Java (JDK) e compilação do servidor \textit{PaperMC}, configurando scripts de inicialização otimizados com parâmetros de memória da JVM recomendados pela comunidade técnica (\textit{Aikar's flags});
    \item Integração dos módulos \textit{GeyserMC} e \textit{Floodgate}, validando a comunicação híbrida entre jogadores que utilizam computadores e dispositivos móveis/consoles.
\end{enumerate}

\subsection{Etapa 3: Aquisição de Domínio, Delegação de DNS e Camada de Segurança}
Visando eliminar a fragilidade de domínios públicos de terceiros e elevar o padrão de resiliência, a metodologia evoluiu para uma arquitetura com domínio próprio:
\begin{enumerate}
    \item Registro de domínio próprio junto ao \textit{Registro.br} e alteração dos servidores de nome autoritativos (\textit{nameservers}) para a rede da \textit{Cloudflare};
    \item Configuração de registros DNS (\textit{A} e \textit{SRV}) para roteamento adequado das conexões de entrada do jogo;
    \item Ativação dos recursos de mitigação de incidentes da Cloudflare, protegendo o ponto de origem contra sobrecargas e ataques de negação de serviço.
\end{enumerate}

\subsection{Etapa 4: Fortalecimento da Infraestrutura, Túnel SSH e Filtragem de Tráfego}
Para consolidar um modelo de administração remota seguro (\textit{hardening}), implementaram-se defesas em profundidade:
\begin{enumerate}
    \item Implantação do \textit{Cloudflare Tunnel} (\textit{cloudflared}), estabelecendo uma conexão criptografada de saída da VPS diretamente para a rede de borda da Cloudflare. Esse mecanismo permitiu fechar integralmente a porta 22 de entrada nas \textit{Security Lists} e no \textit{iptables}, extinguindo qualquer redirecionamento público de portas para fins administrativos e mantendo o acesso SSH viável exclusivamente por meio do túnel autenticado;
    \item Criação de script de rotina, executado periodicamente para consulta de blocos CIDR alocados regionalmente (utilizando listas de endereçamento nacional fornecidas por bases de dados públicas como o \textit{ipdeny.com}), aplicando regras no \textit{firewall} para bloquear tráfego automatizado e conexões maliciosas provenientes de fora da zona de operação pretendida;
    \item Configuração do servidor \textit{Nginx} para publicação segura de serviços complementares (como o mapa tridimensional gerado pelo \textit{BlueMap}), com diretivas para posterior rastreamento e indexação em mecanismos de busca na web.
\end{enumerate}
\chapter{RESULTADOS E DISCUSSÕES}
\label{capitulo-resultados-discussoes}

    Ainda não tem nenhum resultado e discussão, pois o projeto ainda não foi concluído, marco necessário para confecção deste capítulo.
    
	Os resultados podem ser apresentados em forma de tabelas ou gráficos, sendo numerados sequencialmente e discutidos antes de serem colocados. Quando possível os resultados experimentais obtidos devem ser comparados com dados de literatura e suas diferenças (quando houver) discutidas.
	
	Compõe-se de: escrita do assunto a ser tratado no relatório, relato do fato, informação que se queira, descrição de norma e do procedimento. Apresentar somente assuntos que tenham relação direta e específica com o trabalho. 	Interprete as observações e os resultados obtidos. Procure estabelecer relação entre causa e efeito, partindo de suas observações. Mencione os problemas encontrados durante as atividades do projeto e procure relacionar os resultados obtidos com essas dificuldades. A discussão é a sua interpretação das observações e dos resultados. 
	
	Aqui o aluno deve apresentar os resultados esperados e pode também apresentar resultados semelhantes encontrados na literatura em estudos semelhantes realizados por outros autores que já pesquisaram a mesma área objeto da ação de extensão. 
	
	Caso tenha resultados prévios pode apresentar, mas nestes deve-se evitar a repetição por escrito de dados apresentados em tabelas e gráficos e as legendas desses devem ser autoexplicativas, dispensado a leitura do texto para compreensão. É muito importante explicar o que as respostas obtidas significam e como elas podem ajudar a solucionar o problema apresentado na extensão.

\chapter{CONSIDERAÇÕES FINAIS E SUGESTÕES}
\label{capitulo-consideracoes-finais}
	Esta é a parte final do Relatório, na qual o aluno deve apresentar as principais conclusões alcançadas com a atividade desenvolvida.

	Deve expor como a atividade foi importante para sua formação técnica profissional e ressaltar de que maneira os conhecimentos práticos advindos da atividade desenvolvida contribuirão para o conhecimento teórico a ser obtido no curso.

	Com base no trabalho corrente, poderão ser indicados assuntos de relevância significativa para serem abordados em projetos posteriores (sugestões, próximos passos, melhorias, mudanças no encaminhamento das atividades, contribuições do trabalho para a vida acadêmica dos estudantes proponentes da extensão). 
	
    \section{Uso de Inteligência Artificial Generativa}
    \label{capitulo-consideracoes-iag}
    A Portaria CNPq n\textordmasculine  2.664/2026, publicada em 6 de março de 2026, institui a Política de Integridade na Atividade Científica e estabelece as primeiras diretrizes nacionais para o uso de Inteligência Artificial Generativa (IAG) em pesquisas apoiadas pelo CNPq.

    Neste sentido, o uso de Inteligência Artificial Generativa (IAG) pelos alunos será permitida, desde que os alunos informem o seu uso de acordo com os termos da portaria. Em todas as atividades de entrega de trabalhos, os alunos deverão informar os prompts de pesquisa utilizados, qual a IAG utilizada e como ela ajudou no trabalho. Caso seja detectado uso de IAG nas atividades entregues e o aluno não tenha seguido essas regras, o trabalho será considerado inválido e não contará para fins de avaliação. Veja o apêndice \ref{apendice_declaracao_iag} para ver alguns modelos de declaração.

    Outros possíveis declarações para uso de IAG:
    \begin{itemize}
        \item Declaração de uso de IAG para geração/apoio na análise de dados;
        \item Declaração de uso de IAG para geração de código/scripts de pesquisa;
        \item Declaração de uso de IAG para tradução de texto;
        \item Declaração de uso de IAG para geração ou edição de imagens, gráficos e figuras;
        \item Declaração de uso de IAG para transcrição de entrevistas/áudios;
        \item Declaração de uso de IAG para elaboração de resumos/abstracts;
        \item Declaração de uso de IAG na concepção do projeto de pesquisa (brainstorming, formulação de hipóteses);
        \item Declaração de uso de IAG na elaboração de instrumentos de coleta (questionários, roteiros de entrevista);
        \item Declaração de não utilização de IAG (declaração negativa, quando nenhuma ferramenta foi usada);
        \item Declaração de uso de IAG por membros de comitês/pareceristas (uso restrito, vedada inserção de projetos de terceiros);
        \item Declaração de uso de IAG na etapa de submissão (formatação, adequação a normas editoriais).
    \end{itemize}

    Então adapte os modelos sugeridos e faça uma declaração para cada uso de IAG que você.	


%este capítulo não é obrigatório, por isso está no meio do texto. Ele serve como apoio e orientação do documento.
\chapter{REFERÊNCIAS BIBLIOGRÁFICAS}
\label{referencias-bibliograficas}
	Todos os materiais consultados e/ou citados, devem vir nas referências . Para isso, siga as normas contidas no Manual de normas para apresentação de trabalhos acadêmicos do Instituto Federal do Paraná.

	É importante anotar e listar aqui todos os livros, sites, artigos, vídeos, que você referenciou ao longo do trabalho. Evite fazer isso ao final, assim a chance de esquecer alguma referência aumenta.

	As referências dos documentos citados no texto devem ser apresentadas em ordem alfabética, ao final do trabalho.

    Outras dúvidas sobre as normas técnicas sobre as normas técnicas de entrega de trabalhos no IFPR podem ser encontradas no manual \href{https://editora.ifpr.edu.br/index.php/aeditora/catalog/book/72}{Manual de Normas para Apresentação de Trabalhos Acadêmicos do Instituto Federal do Paraná: Padrão ABNT}.

	Preocupe-se em montar o arquivo bibtex do \LaTeX com cuidado e utilizar o comando \verb|\cite| com cuidado durante o texto, que as referências serão montadas automaticamente e dentro das normas.

%este capítulo não é obrigatório, por isso está no meio do texto. Ele serve como apoio e orientação do documento.	
\chapter{AJUDA COM \LaTeX}
	Alguns links para facilitar a vida com o \LaTeX, além de sempre poder utilizar a própria Inteligência Artificial para auxiliar com algumas coisas:

	\begin{itemize}
		\item Link para este documento no Overleaf: \url{https://www.overleaf.com/read/mscjtgrycfvf#063e2d}
		\item Link para meu GitHub (com alguns exemplos de tabelas, apresentação, figuras, algoritmos e link para livro PDF que ensina \LaTeX:) \url{https://github.com/mlaureano/LaTeX}
		\item Muitos exemplos de gráficos \url{https://texample.net/};
		\item Equivalente ao Stack Overflow para programação, só que para o \LaTeX  \url{https://tex.stackexchange.com/}
		\item Tutoriais do Overleaf \url{https://docs.overleaf.com/getting-started/latex-tutorials};
		\item Para ajudar a criar tabelas no \LaTeX: \url{https://www.tablesgenerator.com/}, sempre configure para o estilo ``Bookstabs table style'';
		\item Para ajudar na hora de criar equações matemáticas no \LaTeX: \url{https://editor.codecogs.com/};
		\item Grupo de usuários do \LaTeX no Brasil: \href{https://groups.google.com/g/latex-br?pli=1}{Latex-BR};
		\item Vários exemplos de como fazer gráficos com pgfplots: \url{https://pgfplots.sourceforge.net/gallery.html}
        \item Referência rápida para o \LaTeX: \url{https://quickref.me/latex.html}.
	\end{itemize}
	


% ----------------------------

% ----------------------------------------------------------
% ELEMENTOS PÓS-TEXTUAIS
% ----------------------------------------------------------
\postextual
% ----------------------------------------------------------

% ----------------------------------------------------------
% Referências bibliográficas OBRIGATORIO
% ----------------------------------------------------------
\bibliography{referencias/referencias}

% ----------------------------------------------------------
% Glossário
% ----------------------------------------------------------
%
% Consulte o manual da classe abntex2 para orientações sobre o glossário.
%
%\glossary

% ----------------------------------------------------------
% Apêndices
% ----------------------------------------------------------

% ---
% Inicia os apêndices OPCIONAL
% ---
\begin{apendicesenv}

% Imprime uma página indicando o início dos apêndices
\partapendices

% ----------------------------------------------------------

% =====================================================================
% APÊNDICE — DECLARAÇÕES DE USO DE INTELIGÊNCIA ARTIFICIAL GENERATIVA (IAG)
% Conforme Portaria CNPq nº 2.664, de 6 de março de 2026
% (Institui a Política de Integridade na Atividade Científica do CNPq)
%
% Este arquivo foi pensado para ser incluído via \input{} ou \include{}
% em um relatório maior, dentro do ambiente de apêndices.
% Requer os pacotes: geometry (opcional), enumitem, array, xcolor (opcional)
% =====================================================================

\chapter{Declarações de Uso de Inteligência Artificial Generativa (IAG)}
\label{apendice_declaracao_iag}

\noindent
Em conformidade com a Portaria CNPq nº 2.664, de 6 de março de 2026, que institui a
Política de Integridade na Atividade Científica do CNPq e, em seu Art.~9º, inciso I,
alíneas ``c'' e ``d'', determina a obrigatoriedade de declaração do uso de ferramentas
de Inteligência Artificial Generativa (IAG) em qualquer fase do desenvolvimento da
pesquisa --- concepção, redação, análise de dados ou submissão ---, especificando a
ferramenta utilizada e a finalidade, apresentam-se a seguir as declarações relativas
ao uso de IAG neste trabalho.

\bigskip
\noindent
Ressalta-se que, nos termos da referida Portaria, é vedada a submissão de conteúdo
gerado por IAG como se fosse de autoria humana, sendo os autores integralmente
responsáveis pelo conteúdo final, inclusive por eventuais plágios ou imprecisões
geradas pela ferramenta (Art.~9º, I, ``d'' e ``f'').

% =====================================================================
% MODELO 1 — DECLARAÇÃO DE USO DE IAG PARA REVISÃO DE TEXTO
% =====================================================================

\section{Modelo de Declaração --- Revisão de Texto}
\label{ap:declaracao_revisao_texto}

\begin{center}
\textbf{\large DECLARAÇÃO DE USO DE INTELIGÊNCIA ARTIFICIAL GENERATIVA}\\
\textbf{(Revisão Textual e Linguística)}
\end{center}

\bigskip
\noindent
Eu, \underline{\hspace{7cm}} (nome completo do(a) autor(a) ou responsável), aluno da disciplina Práticas de Extensão I, II, III ou IV, DECLARO, para os devidos fins e em atendimento ao disposto na Portaria CNPq nº 2.664/2026, que fiz uso de ferramenta(s) de Inteligência Artificial Generativa (IAG) no desenvolvimento deste trabalho, conforme especificado abaixo:

\bigskip
\begin{tabular}{|p{4.5cm}|p{9.5cm}|}
\hline
\textbf{Ferramenta(s) utilizada(s)} & \underline{\hspace{9cm}} \\
& (ex.: ChatGPT, Claude, Grammarly, DeepL Write, etc. --- indicar nome e versão/modelo) \\
\hline
\textbf{Fase(s) da pesquisa em que houve uso} &
\begin{itemize}[leftmargin=*, nosep]
  \item[$\square$] Concepção
  \item[$\square$] Redação
  \item[$\square$] Análise de dados
  \item[$\square$] Submissão
\end{itemize} \\
\hline
\textbf{Finalidade do uso} & Revisão gramatical, ortográfica e de estilo do texto;
melhoria de clareza e coesão textual; correção de fluência da linguagem.
\emph{(Especificar detalhes adicionais, se houver)}: \underline{\hspace{9cm}} \\
\hline
\textbf{Trechos/seções do trabalho em que a ferramenta foi aplicada} &
\underline{\hspace{9cm}} \\
\hline
\end{tabular}

\bigskip
\noindent
Declaro, ainda, que:
\begin{enumerate}[label=\alph*)]
  \item o conteúdo intelectual, as análises, as interpretações e as conclusões
  apresentadas neste trabalho são de minha inteira autoria e responsabilidade;
  \item a ferramenta de IAG foi empregada exclusivamente como apoio à revisão
  linguística e estilística do texto, não tendo sido utilizada para geração de
  conteúdo científico, dados, resultados ou conclusões;
  \item todo o texto revisado por IAG foi integralmente analisado, verificado e
  validado por mim antes da submissão, assumindo integral responsabilidade por
  eventuais imprecisões, distorções de sentido ou plágios porventura introduzidos
  pela ferramenta;
  \item não foi submetido conteúdo gerado por IAG como se fosse de autoria humana,
  em conformidade com o Art.~9º, I, ``d'', da Portaria CNPq nº 2.664/2026.
\end{enumerate}

\bigskip
\noindent
Pinhais e data: \underline{\hspace{4cm}}, \underline{\hspace{1.5cm}} de
\underline{\hspace{3cm}} de \underline{\hspace{1.5cm}}.

\bigskip\bigskip
\noindent
\underline{\hspace{7cm}}\\
Assinatura do(a) autor(a) / responsável

% =====================================================================
% MODELO 2 — DECLARAÇÃO DE USO DE IAG PARA REVISÃO DE BIBLIOGRAFIA
% =====================================================================

\newpage
\section{Modelo de Declaração --- Revisão e Levantamento Bibliográfico}
\label{ap:declaracao_revisao_bibliografia}

\begin{center}
\textbf{\large DECLARAÇÃO DE USO DE INTELIGÊNCIA ARTIFICIAL GENERATIVA}\\
\textbf{(Levantamento, Revisão e Organização Bibliográfica)}
\end{center}

\bigskip
\noindent
Eu, \underline{\hspace{7cm}} (nome completo do(a) autor(a) ou responsável),
aluno da disciplina Práticas de Extensão I, II, III ou IV, DECLARO, para os devidos
fins e em atendimento ao disposto na Portaria CNPq nº 2.664/2026, que fiz uso de
ferramenta(s) de Inteligência Artificial Generativa (IAG) para apoio ao levantamento,
à revisão e/ou à organização de referências bibliográficas neste trabalho, conforme
especificado abaixo:

\bigskip
\begin{tabular}{|p{4.5cm}|p{9.5cm}|}
\hline
\textbf{Ferramenta(s) utilizada(s)} & \underline{\hspace{9cm}} \\
& (ex.: ChatGPT, Claude, Perplexity, Elicit, Consensus, Scite, etc. --- indicar nome
e versão/modelo) \\
\hline
\textbf{Fase(s) da pesquisa em que houve uso} &
\begin{itemize}[leftmargin=*, nosep]
  \item[$\square$] Concepção
  \item[$\square$] Redação
  \item[$\square$] Análise de dados
  \item[$\square$] Submissão
\end{itemize} \\
\hline
\textbf{Finalidade do uso} &
\begin{itemize}[leftmargin=*, nosep]
  \item[$\square$] Busca e sugestão de referências relacionadas ao tema
  \item[$\square$] Organização/formatação de referências (normas ABNT, APA, Vancouver etc.)
  \item[$\square$] Síntese ou resumo de artigos previamente selecionados pelo(a) autor(a)
  \item[$\square$] Verificação de completude/consistência de citações no texto
  \item[$\square$] Outra finalidade: \underline{\hspace{5cm}}
\end{itemize} \\
\hline
\end{tabular}

\bigskip
\noindent
Declaro, ainda, que:
\begin{enumerate}[label=\alph*)]
  \item todas as referências sugeridas ou organizadas com apoio da ferramenta de IAG
  foram individualmente verificadas quanto à sua existência, exatidão bibliográfica
  e pertinência ao tema, mediante consulta às fontes originais;
  \item nenhuma referência foi incluída no trabalho sem confirmação de sua
  autenticidade, sendo de minha exclusiva responsabilidade eventuais erros,
  omissões ou citações fabricadas (\emph{``alucinações''}) porventura geradas pela
  ferramenta;
  \item a seleção crítica, a interpretação e o uso das fontes bibliográficas no
  corpo do trabalho são de minha inteira autoria e responsabilidade;
  \item não foram inseridos, nas ferramentas de IAG utilizadas, projetos de pesquisa
  de terceiros ou dados de caráter sigiloso para fins de elaboração de pareceres,
  buscas ou análises, em conformidade com o Art.~9º, I, ``e'', da Portaria CNPq
  nº 2.664/2026.
\end{enumerate}

\bigskip
\noindent
Pinhais e data: \underline{\hspace{4cm}}, \underline{\hspace{1.5cm}} de
\underline{\hspace{3cm}} de \underline{\hspace{1.5cm}}.

\bigskip\bigskip
\noindent
\underline{\hspace{7cm}}\\
Assinatura do(a) autor(a) / responsável

% =====================================================================
% Nota de rodapé de embasamento legal (opcional — pode ser removida)
% =====================================================================
\bigskip
\noindent\rule{\textwidth}{0.4pt}
\footnotesize
\noindent
\textbf{Base legal:} Portaria CNPq nº 2.664, de 6 de março de 2026, Art.~9º, inciso I,
alíneas ``c'' (dever de declarar o uso de IAG, especificando ferramenta e finalidade),
``d'' (vedação à submissão de conteúdo de IAG como autoria humana), ``e'' (vedação à
inserção de projetos de pesquisa de terceiros em ferramentas de IAG) e ``f''
(responsabilidade integral do autor pelo conteúdo final, inclusive por plágios ou
imprecisões geradas pela IAG).
\normalsize




%se tiver mais apendices, crie um novo arquivo tex e coloque no diretório.
% ----------------------------------------------------------
\chapter{Nullam elementum urna vel imperdiet sodales elit ipsum pharetra ligula
ac pretium ante justo a nulla curabitur tristique arcu eu metus}
% ----------------------------------------------------------
\lipsum[55-57]

\end{apendicesenv}
% ---


% ----------------------------------------------------------
% Anexos
% ----------------------------------------------------------

% ---
% Inicia os anexos OPCIONAL
% ---
\begin{anexosenv}

% Imprime uma página indicando o início dos anexos
\partanexos
%% crie um arquivo .tex para cada anexo.
% ---
\chapter{Morbi ultrices rutrum lorem.}
% ---
\lipsum[30]

% ---
\chapter{Cras non urna sed feugiat cum sociis natoque penatibus et magnis dis
parturient montes nascetur ridiculus mus}
% ---

\lipsum[31]

% ---
\chapter{Fusce facilisis lacinia dui}
% ---

\lipsum[32]

\end{anexosenv}



\end{document}
